\documentclass[../../main.tex]{subfiles}

\graphicspath{{../../image/}{../../figure/}}

\begin{document}

\section{注记}\label{section:ch01-注记}

平衡律理论的原理是在十九世纪二十年代奠定弹性力学基础的过程中形成的。\cref{theorem:ch01-场方程归约} 有悠久而著名的历史。通量密度必为外法向线性函数这一关键发现由 Cauchy [1,2] 作出；通量密度通过曲面仅依赖于外法向这一论证归功于 Hamel 与 Noll [2]。关于这些思想在连续介质物理背景下的近期发展，见 Gurtin and Martins [1]、Degiovanni, Marzocchi and Musesti [1]、Marzocchi and Musesti [1,2] 与 Šilhavý [2,3,4]。此处证明借自 Ziemer [1]。关于从场方程恢复平衡律这一问题——由\cref{theorem:ch01-一般迹} 处理——Chen 与 Frid [1,5,6] 发展了一套完整的散度测度场理论，在曲面可叶状化这一附加温和技术性假设下（类似于\cref{remark:ch01-显式迹构造}），给出了比本文方法更显式的迹构造。该途径的进一步发展见 Chen [9,10]、Chen and Frid [8,9]、Frid [7]、Chen and Torres [1,2] 以及 Chen, Torres and Ziemer [1,2]。另一种不那么显式、泛函分析风格的途径见 Anzellotti [1]。一个正在研究中的重要问题是：当坐标变换属于某 Sobolev 空间 $W^{1,p}$ 甚至 $BV$ 时，\cref{theorem:ch01-坐标变换} 的结论是否仍成立。

平衡律方程组配有非平凡伴随当且仅当它可对称化这一观察归功于 Godunov [1,2,3] 以及 Friedrichs and Lax [1]；另见 Boillat [1] 与 Ruggeri and Strumia [1]。关于物理中出现的平衡律方程组正则截断的讨论，见 Boillat and Ruggeri [1]。

如历史引论中已述，在一维空间中，弱阵面最早出现于 Euler 的声学研究，而激波阵面由 Stokes [1] 引入。多维空间中的阵面最早由 Christoffel [1] 研究。激波阵面与相变之间的联系此处不作探讨。这一活跃研究领域的参考文献见%

\cref{chapter:ch08-容许激波}。

$BV$ 函数理论的系统论述可见 Federer 的专著 [1]、Giusti 与 Ambrosio, Fusco and Pallara 的著作 [1]，以及 Evans and Gariepy [1] 与 Ziemer [2]。\cref{theorem:ch01-BV链式法则,theorem:ch01-Gauss-Green} 取自 Volpert [1]。特殊有界变差函数理论在 Ambrosio, Fusco and Pallara [1] 中得到精细阐述。

关于\cref{section:ch01-快速振荡} 所提问题的富有洞见的讨论见 DiPerna [10]。这些问题将在%

\cref{chapter:ch16-平衡律组bv解} 介绍补偿列紧方法时阐明。

\end{document}
